\section{Discussion}
\label{sec:discussion}

\subsection{AI assistance in the development process}
\label{sec:ai}

We used LLM-based coding assistants (Claude Code, Anthropic) to accelerate
formalization, in the same general direction as recent work on agentic
autoformalization~\cite{ren2026merlean}.

Our workflow followed a \emph{sorry-driven development} pattern:
\begin{enumerate}
  \item State all theorem signatures with \texttt{sorry} placeholders.
  \item Verify that the signatures compose correctly (the main theorem
    type-checks with all intermediates as \texttt{sorry}).
  \item Fill \texttt{sorry}s bottom-up, from leaves of the dependency DAG
    to the root.
\end{enumerate}
This approach worked well with LLM agents: independent leaves
(e.g., the four exponential bounds B1--B4) could be assigned to parallel
agents, each given the exact signature to fill and the list of available
lemmas.

\textbf{Where agents excelled.} ``Fill this \texttt{sorry} given these
lemmas''---tasks with a clear specification and known
ingredients---compiled on the first attempt when delegated to agents.
Most of the exponential bounds and the assembly step were completed this way.

\textbf{Where agents struggled.} ``Find the right proof approach''---tasks
requiring mathematical insight, such as the Strang commutator-scaling argument
or the algebraic factorization for the step error---required human-driven
exploration. The agent's tendency to attempt minor variations of a failing
approach (rather than reconsidering the overall strategy) became a bottleneck
for these tasks.

\textbf{A cautionary pattern.} As noted by Zhao et al.~\cite{zhao2025chsh},
LLM agents sometimes quietly strengthen hypotheses to make a proof trivial.
The \texttt{sorry}-first workflow guards against this: since all signatures are
fixed before proof filling begins, adding hypotheses breaks the type-checking
of downstream theorems.

\subsection{Mathlib ecosystem}
\label{sec:mathlib}

Several gaps in \Mathlib were encountered and filled during this project:
\begin{itemize}
  \item $\norm{e^a} \le e^{\norm{a}}$ for general Banach algebras
    (\leanref{ExpBounds.lean}{72}{norm\_exp\_le}).
  \item $e^x - 1 - x \le x^2/2 \cdot e^x$ for $x \ge 0$
    (\leanref{ExpBounds.lean}{129}{exp\_sub\_one\_sub\_bound\_real}).
  \item $\norm{e^a - 1} \le e^{\norm{a}} - 1$
    (\leanref{ExpBounds.lean}{88}{norm\_exp\_sub\_one\_le}).
\end{itemize}
These are natural candidates for upstreaming to \Mathlib.

Working over a general field $\mathbb{K}$ (via \texttt{RCLike}) introduced
friction with the calculus infrastructure, which assumes $\mathbb{R}$-module
structure. Our pragmatic solution---working over $\mathbb{R}$ directly in the
commutator-scaling files---is effective but means users must apply
\texttt{restrictScalars} for $\mathbb{C}$ applications. A unified treatment
would require additional typeclass instances in \Mathlib.

\subsection{Formalization-driven proof design}
\label{sec:formalization-insights}

The formalization process influenced proof design in several ways:

\textbf{Algebraic factorization.}
The standard approach to the step error---expanding $e^a\,e^b$ to second order
and bounding cross-terms---requires tracking non-commutative multinomial
coefficients. In Lean, this involves dozens of \texttt{norm\_mul\_le} and
\texttt{gcongr} steps for each cross-term. The algebraic factorization
(\cref{sec:step-error}) splits the problem into two independent bounds,
each handled by series comparison, reducing the proof length by approximately
half. This factorization is arguably cleaner than the standard approach even
on paper.

\textbf{Duhamel without Gronwall.}
The commutator-scaling proof was designed with Lean feasibility in mind.
The standard approach uses Gronwall's inequality or ODE uniqueness, which
would require formalizing existence and uniqueness for linear ODEs in
Banach spaces---theory not yet available in \Mathlib. The Duhamel/FTC approach
uses only the product rule and FTC-2, both already available.

\textbf{\texttt{noncomm\_ring} vs.\ \texttt{ring}.}
The \texttt{ring} tactic in Lean silently fails on non-commutative goals:
it produces an unsolved subgoal rather than an error message. We discovered
this only after extended debugging. The \texttt{noncomm\_ring} tactic handles
free non-commutative ring identities correctly but cannot use commutativity
hypotheses (e.g., $[a, e^a] = 0$). Our workaround: use \texttt{set} to make
exponentials opaque, rewrite commutativity as equalities, then apply
\texttt{noncomm\_ring} on the resulting free-ring identity.

\subsection{Relation to quantum simulation}
\label{sec:quantum}

Our results are stated for general complete normed algebras, which include
the bounded operators $B(\mathcal{H})$ on a Hilbert space as a special case.
Specialization to finite-dimensional matrix algebras
$M_d(\mathbb{C})$ is straightforward via \Mathlib's
\texttt{Matrix.instNormedRing}.

The commutator-scaling bounds are physically motivated by Hamiltonian
simulation: for a lattice Hamiltonian $H = \sum_j h_j$ where each $h_j$ acts
locally, the commutator $\comm{h_j}{h_k}$ vanishes for distant terms.
Our formalized bounds could serve as certified building blocks for verified
resource estimation in quantum algorithms, once combined with gate-count
analysis for individual exponentials.
